\section{Quickstart}
%%%%%%%%%%%%%%%%%%%%%%%%%%%%%%%%%%%%%%%%%%%%%%%%%%%%%%%%%%%%%%%%%%%%%%%%%%%%%%%%
\subsection{What is \taoce}
\begin{frame}
    \taoce is an open-source platform developed by \OATSA for online testing and assessment.
    \only<article>{%
        \TAOCE has been released in january 2026, and a new version has been released in september.

        The platform is published as container images, and can run on different systems, like laptop, virtual machines or even Raspberry Pi.

        The following applications support core features of \taoce
    }
\pause
    \begin{description}
        \item[Portal] Central gateway of \taoce, to manage users, groups, sessions, \dots
        \pause
        \item[Deliver] (aka. Testrunner) is where test-takers will attend their sessions.
        \pause
        \item[Proctoring] (aka. Monitoring) is a parallel service to monitor and control test sessions.
        \pause
        \item[Scoring] (aka. Grader) can be used to support manual scoring sessions.
        \pause
        \item[Construct] (aka. Authoring) is the historical platform for TAO Core. In \taoce, it offers an authoring interface to create QTI items and tests. 
    \end{description}
\end{frame}
%%%%%%%%%%%%%%%%%%%%%%%%%%%%%%%%%%%%%%%%%%%%%%%%%%%%%%%%%%%%%%%%%%%%%%%%%%%%%%%%
\subsection{\TAOCE Applications \& Flavors}
\newcommand{\flyes}{\alert{\faCheckCircle}}
\begin{frame}[fragile]
    \only<article>{%
        \taoce come in different set of applications, called \alert{flavors} to match different use-cases.
    }
    \begin{center}
        \begin{tabular}{l | c c c}
            & \scriptsize\texttt{full}
            & \scriptsize\texttt{essential}
            & \scriptsize\texttt{lite}
            \\
            \hline
            Proctoring      & \flyes & & \\
            Scoring         & \flyes & & \\
            Construct       & \flyes & \flyes \\
            LTI Devkit      & \flyes & \flyes \\
            Portal          & \flyes & \flyes & \flyes \\
            Deliver         & \flyes & \flyes & \flyes \\
            EM \& internal  & \flyes & \flyes & \flyes \\
        \end{tabular}

        \vspace*{2em}
        We will focus on the \alert{lite} flavor for this presentation.
    \end{center}
\end{frame}
%%%%%%%%%%%%%%%%%%%%%%%%%%%%%%%%%%%%%%%%%%%%%%%%%%%%%%%%%%%%%%%%%%%%%%%%%%%%%%%%
\subsection{How we design \taoce}
\begin{frame}
    \taoce is composed in three layers:

    \begin{columns}[T]
        \begin{column}{.3\textwidth}\begin{center}
            \vspace{1em}
            \begin{figure}[!htb]
                \begin{tikzpicture}[scale=2]
                    \draw<1->[thick,tao-blue!90] (0.667,0.577) -- (1.333,0.577) -- (1,1.155) -- cycle;
                    \draw<2->[thick,tao-blue!70] (0.333,0.289) -- (1.667,0.289) -- (1,1.443) -- cycle;
                    \draw<3->[thick,tao-blue!50] (0,0) -- (2,0) -- (1,1.732) -- cycle;
                    \draw[thick,transparent] (0,0) -- (2,0) -- (1,1.732) -- cycle;

                    \node<1->[font=\scriptsize] at (1, 0.666) {Crystal};
                    \node<2->[font=\scriptsize] at (1, 0.433) {Swift};
                    \node<3->[font=\scriptsize] at (1, 0.145) {Cozy};
                \end{tikzpicture}
            \end{figure}
        \end{center}\end{column}
        \begin{column}{.7\textwidth}
            \begin{description}
                \item<1->[Crystal] the set of \taoce applications built from sources
                \item<2->[Swift] the packaging of \alert{Crystal} as container images
                \item<3->[Cozy\notcovered] an appliance to run container from \alert{Swift} in a Virtual Machine or on Bare Metal
            \end{description}
        \end{column}
    \end{columns}
\end{frame}
%%%%%%%%%%%%%%%%%%%%%%%%%%%%%%%%%%%%%%%%%%%%%%%%%%%%%%%%%%%%%%%%%%%%%%%%%%%%%%%%
\subsection{Start Development Environment}
\begin{frame}[fragile]
    \only<article>{%
        \taoce is built in an isolated container environment, called \devcontainer\notcovered.
    }
    % \bigicon{\faPlayCircle}
    \begin{description}
        \item[Get code] \,
            \begin{lstlisting}
git clone \
  --recurse-submodules \
  git@github.com:tao-ce/tao-ce.git
\end{lstlisting}
        \item[Enter \devcontainer\notcovered] \, \\
            In Visual Studio Code \\
            \keys{\ctrl + \shift + P} Dev Containers: Reopen in Container
        \item[Build \& Run] \, \\
            \begin{lstlisting}
# Limit scope to lite flavor
export CE_FLAVOR=lite
# Build & run TAO CE from devcontainer
task dev:up
\end{lstlisting}
    \end{description}
\end{frame}
%%%%%%%%%%%%%%%%%%%%%%%%%%%%%%%%%%%%%%%%%%%%%%%%%%%%%%%%%%%%%%%%%%%%%%%%%%%%%%%%
% \begin{frame}
% After few minutes, a container with the complete toolchain to build \taoce Applications and Images will be ready.\\

% \bigicon{\faHourglass}

% We will not cover in details how Devcontainers are working. For more information:
% \begin{itemize}
%     \item \href{https://containers.dev/}{Development Containers website}
%     \item \href{https://github.com/tao-ce/tao-ce/blob/main/.devcontainer/devcontainer.json}{\TAOCE \devcontainer}
% \end{itemize}

% \end{frame}